The following functions are intended as a way for API users to be
informed when an error or significant event occurs.
Each function allows a user to register a handler for an event. The
return code for all callback registration
functions is the address of the handler that was previously
registered (which may be NULL if no handler was previously
registered). For backwards compatibility reasons, some callbacks may
pass a \code{BPatch\_thread} object when a
\code{BPatch\_process} may be more appropriate. A
\code{BPatch\_thread} may be converted into a \code{BPatch\_process} using
\code{BPatch\_thread::getProcess()}.

\subsection{Asynchronous Callbacks}

\begin{apient}
  typedef void (*BPatchAsyncThreadEventCallback)(
    BPatch_process *proc,
    BPatch_thread *thread
  )

  bool registerThreadEventCallback(BPatch_asyncEventType type,
                                   BPatchAsyncThreadEventCallback cb)

  bool removeThreadEventCallback(BPatch_asyncEventType type,
                                 BPatch_AsyncThreadEventCallback cb)
\end{apient}

\apidesc{
  The type parameter can be either one of
  \code{BPatch\_threadCreateEvent} or
  \code{BPatch\_threadDestroyEvent}. Different callbacks
  can be registered for different values of type.
}

\subsection{Code Discovery Callbacks}

\begin{apient}
  typedef void (*BPatchCodeDiscoveryCallback)(
    BPatch_Vector<BPatch_function*> &newFuncs,
    BPatch_Vector<BPatch_function*> &modFuncs
  )

  bool registerCodeDiscoveryCallback(BPatchCodeDiscoveryCallback cb)
  bool removeCodeDiscoveryCallback(BPatchCodeDiscoveryCallback cb)
\end{apient}

\apidesc{
  This callback is invoked whenever previously un-analyzed code is
  discovered through runtime analysis, and delivers a
  vector of functions whose analysis have been modified and a vector
  of functions that are newly discovered.
}

\subsection{Code Overwrite Callbacks}

\begin{apient}
  typedef void (*BPatchCodeOverwriteBeginCallback)(
    BPatch_Vector<BPatch_basicBlock*> &overwriteLoopBlocks
  )

  typedef void (*BPatchCodeOverwriteEndCallback)(
    BPatch_Vector<std::pair<Dyninst::Address,int>> &deadBlocks,
    BPatch_Vector<BPatch_function*> &owFuncs,
    BPatch_Vector<BPatch_function*> &modFuncs,
    BPatch_Vector<BPatch_function*> &newFuncs
  )

  bool registerCodeOverwriteCallbacks(BPatchCodeOverwriteBeginCallback cbBegin,
                                      BPatchCodeOverwriteEndCallback cbEnd)
\end{apient}

\apidesc{
  Register a callback at the beginning and end of overwrite events.
  Only invoke if Dyninst\'s hybrid analysis mode is set
  to \code{BPatch\_defensiveMode}.

  The \code{BPatchCodeOverwriteBeginCallback} callback allows the
  user to remove any instrumentation when the program starts
  writing to a code page, which may be desirable as instrumentation
  cannot be removed during the overwrite loop\'s
  execution, and any breakpoint instrumentation will dramatically
  slow the loop\'s execution.

  The \code{BPatchCodeOverwriteEndCallback} callback delivers the
  effects of the overwrite loop when it is done executing. In
  many cases no code will have changed.
}

\subsection{Dynamic calls}

\begin{apient}
  typedef void (*BPatchDynamicCallSiteCallback)(
    BPatch_point*at_point,
    BPatch_function *called_function
  )

  bool registerDynamicCallCallback(BPatchDynamicCallSiteCallback cb)
  bool removeDynamicCallCallback(BPatchDynamicCallSiteCallback cb)
\end{apient}

\apidesc{
  The \code{registerDynamicCallCallback} interface will not
  automatically instrument any dynamic call site. To make sure the call
  back function is called, the user needs to explicitly instrument
  dynamic call sites. One way to achieve this goal is to
  first get instrumentation points representing dynamic call sites
  and then call \code{BPatch\_point::monitorCalls} with a NULL
  input parameter.
}

\subsection{Dynamic libraries}

\begin{apient}
  typedef void (*BPatchDynLibraryCallback)(
    BPatch_thread *thr,
    BPatch_object *obj,
    bool loaded
  )

  BPatchDynLibraryCallback registerDynLibraryCallback(BPatchDynLibraryCallback func)
\end{apient}

\apidesc{
  Note that in versions previous to 9.1, BPatchDynLibraryCallback\'s
  signature took a \code{BPatch\_module}
  instead of a \code{BPatch\_object}.
}

\subsection{Errors}

\begin{apient}
  enum BPatchErrorLevel {
    BPatchFatal,
    BPatchSerious,
    BPatchWarning,
    BPatchInfo
  }
\end{apient}

\begin{apient}
  typedef void (*BPatchErrorCallback)(
    BPatchErrorLevel severity,
    int number,
    const char * const *params
  )
\end{apient}

\begin{apient}
  BPatchErrorCallback registerErrorCallback(BPatchErrorCallback func)
\end{apient}

\apidesc{
  This function registers the error callback function with the BPatch
  class. The return value is the address of the
  previous error callback function. Dyninst users can change the
  error callback during program execution (e.g., one error
  callback before a GUI is initialized, and a different one after).
  The severity field indicates how important the error
  is (from fatal to information/status). The number is a unique
  number that identifies this error message. Params are
  the parameters that describe the detail about an error, e.g., the
  process id where the error occurred. The number and
  meaning of params depends on the error. However, for a given error
  number the number of parameters returned will
  always be the same.
}

\subsection{Exec}

\begin{apient}
  typedef void (*BPatchExecCallback)(
    BPatch_thread *thr
  )
\end{apient}

\begin{apient}
  BPatchExecCallback registerExecCallback(BPatchExecCallback func)
\end{apient}

\apidesc{
  Not implemented on Windows.
}

\subsection{Exit}

\begin{apient}
  enum BPatch_exitType {
    NoExit,
    ExitedNormally,
    ExitedViaSignal
  }
\end{apient}

\apidesc{
  Mechanism by which the process exited.
}

\begin{apient}
  typedef void (*BPatchExitCallback)(
    BPatch_thread *proc,
    BPatch_exitType exit_type
  )
\end{apient}

\begin{apient}
  BPatchExitCallback registerExitCallback(BPatchExitCallback func)
\end{apient}
  
\apidesc{
  Register a function to be called when a process terminates. For a
  normal process exit, the callback will actually be
  called just before the process exits, but while its process state
  still exists. This allows final actions to be taken
  on the process before it actually exits. The function
  \code{BPatch\_thread::isTerminated()} will return true in this context
  even though the process hasn't yet actually exited. In the case
  of an exit due to a signal, the
  process will have already exited.
}

\subsection{Fork}

\begin{apient}
  typedef void (*BPatchForkCallback)(
    BPatch_thread *parent,
    BPatch_thread *child
  )
\end{apient}

\apidesc{
  This is the prototype for the pre-fork and post-fork callbacks.
  The parent parameter is the parent thread, and the
  child parameter is a \code{BPatch\_thread} in the newly created
  process. When invoked as a pre-fork callback, the child is NULL.
}

\begin{apient}
  BPatchForkCallback registerPreForkCallback(BPatchForkCallback func)
\end{apient}

\apidesc{
  Not implemented on Windows
}

\begin{apient}
  BPatchForkCallback registerPostForkCallback(BPatchForkCallback func)
\end{apient}

\apidesc{
  Not implemented on Windows

  Register callbacks for pre-fork (before the child is created) and
  post-fork (immediately after the child is created).
  When a pre-fork callback is executed the child parameter will be NULL.
}

\subsection{One Time Code}

\begin{apient}
  typedef void (*BPatchOneTimeCodeCallback)(
    Bpatch_thread *thr,
    void *userData,
    void *returnValue
  )
\end{apient}

\begin{apient}
  BPatchOneTimeCodeCallback registerOneTimeCodeCallback(BPatchOneTimeCodeCallback func)
\end{apient}

\apidesc{
  The thr field contains the thread that executed the oneTimeCode
  (if thread-specific) or an unspecified thread in the
  process (if process-wide). The userData field contains the value
  passed to the oneTimeCode call. The returnValue
  field contains the return result of the oneTimeCode snippet.
}

\subsection{Signal Handler}

\begin{apient}
  typedef void (*BPatchSignalHandlerCallback)(
    BPatch_point *at_point,
    long signum,
    std::vector<Dyninst::Address> *handlers
  )
\end{apient}

\begin{apient}
  bool registerSignalHandlerCallback(BPatchSignalHandlerCallback cb, std::set<long> &signal_numbers)
  bool registerSignalHandlerCallback(BPatchSignalHandlerCallback cb, BPatch_Set<long> *signal_numbers)
\end{apient}

\apidesc{
  This function registers the signal handler callback function with
  the BPatch class. The return value indicates success
  or failure. The \code{signal\_numbers} set contains those signal
  numbers for which the callback will be invoked.

  The \code{at\_point} parameter indicates the point at which the
  signal/exception was raised, signum is the number of the
  signal/exception that was raised, and the handlers vector
  contains any registered handler(s) for the signal/exception.
  In Windows this corresponds to the stack of Structured Exception
  Handlers, while for Unix systems there will be at
  most one registered exception handler. This functionality is only
  fully implemented for the Windows platform.
}

\begin{apient}
  bool removeSignalHandlerCallback(BPatchSignalHandlerCallback cb)
\end{apient}

\apidesc{
  Removes a callback registered with \code{registerSignalHandlerCallback}.
}



\subsection{Stopped Threads}

\begin{apient}
  typedef void (*BPatchStopThreadCallback)(
    BPatch_point *at_point,
    void *returnValue
  )
\end{apient}

\apidesc{
  This is the prototype for the callback that is associated with
  the stopThreadExpr snippet class (see Section).
  Unlike the other callbacks in this section, stopThreadExpr
  callbacks are registered during
  the creation of the stopThreadExpr snippet type. Whenever a
  stopThreadExpr snippet executes in a given thread, the
  snippet evaluates the calculation snippet that stopThreadExpr
  takes as a parameter, stops the
  thread\'s execution and invokes this callback. The
  \code{at\_point} parameter is the \code{BPatch\_point} at
  which the stopThreadExpr snippet was inserted, and returnValue
  contains the computation made by the calculation
  snippet.
}

\subsection[User-triggered callbacks]{User-triggered callbacks}

\begin{apient}
  typedef void (*BPatchUserEventCallback)(
    BPatch_process *proc,
    void *buf,
    unsigned int bufsize
  )
\end{apient}

\begin{apient}
  bool registerUserEventCallback(BPatchUserEventCallback cb)
\end{apient}

\apidesc{
  Register a callback that is executed when the user sends a
  message from the mutatee using the DYNINSTuserMessage
  function in the runtime library.
}

\begin{apient}
  bool removeUserEventCallback(BPatchUserEventCallback cb)
\end{apient}

\apidesc{
  Removes a callback registered with \code{registerUserEventCallback}.
}
